% Generated draft: requires mathematical and source review.
\documentclass[11pt]{article}
\usepackage[margin=25mm]{geometry}
\usepackage{fontspec}
\setmainfont{Latin Modern Roman}
\usepackage{amsmath,amssymb,mathtools,unicode-math}
\setmathfont{Latin Modern Math}
\usepackage{enumitem,needspace,xurl,hyperref}
\setlength{\parindent}{0pt}
\setlength{\parskip}{5pt}
\setlength{\emergencystretch}{3em}
\newcommand{\N}{\mathbb{N}}
\newcommand{\Z}{\mathbb{Z}}
\newcommand{\Q}{\mathbb{Q}}
\newcommand{\R}{\mathbb{R}}
\newcommand{\C}{\mathbb{C}}
\newcommand{\F}{\mathbb{F}}
\newcommand{\A}{\mathbb{A}}
\newcommand{\PP}{\mathbb{P}}
\newcommand{\E}{\mathbb{E}}
\newcommand{\eps}{\varepsilon}
\newcommand{\dd}{\,\mathrm{d}}
\begin{document}
\section*{UnsolvedMath ID: 30002004}
\textbf{Source problem code:} catalogue entry
\subsection{Definitions and mathematical statement}
The supplied record does not specify the ground field, the quantum-affine algebra, its deformation parameter, or the category in which extensions are computed. Consequently, the following is a conditional formulation of the recorded characterization problem, pending those choices.

Let \(k\) denote the intended ground field, let \(A\) denote the intended quantum-affine algebra over \(k\), and let \(\mathcal C\) denote the intended \(k\)-linear abelian category of \(A\)-modules. For objects \(V\) of \(\mathcal C\), define \(\operatorname{Ext}^{1}_{\mathcal C}(V,V)\) to be the \(k\)-vector space of Yoneda equivalence classes of short exact sequences in \(\mathcal C\),
\[
0\longrightarrow V\longrightarrow E\longrightarrow V\longrightarrow 0,
\]
where equivalence means an isomorphism of short exact sequences inducing the identity on both end terms, and addition is the Baer sum. A module is simple if it is nonzero and has no nonzero proper submodule.

For the intended fixed choices of \(k\), \(A\), and \(\mathcal C\), characterize, up to \(A\)-module isomorphism, all simple modules \(V\in\mathcal C\) such that
\[
\dim_k V<\infty
\qquad\text{and}\qquad
\dim_k\operatorname{Ext}^{1}_{\mathcal C}(V,V)=1.
\]
The requested characterization should give necessary and sufficient conditions for this equality for every simple finite-dimensional module in the intended category. The unspecified choices above are unresolved qualifications, not an assertion that the same characterization applies to every quantum-affine algebra or every extension category. In particular, the supplied record does not state a precise equivalence between this condition and primeness.
\subsection{Short English statement}
Determine exactly which simple finite-dimensional quantum-affine modules have a one-dimensional space of self-extensions. A precise classification requires specifying the ground field, quantum-affine algebra, deformation parameter, and extension category. The supplied record labels the problem open; this draft does not independently verify that status.
\subsection{Sources}
\begin{itemize}
\item Supplied local UnsolvedMath registry record, permanent ID \(30002004\), title: \emph{Self-Extensions of Simple Quantum-Affine Modules}; source problem code \texttt{catalogue entry}. The record assigns the year 2012 and names no individual proposer.
\item Supplied report provenance: \emph{Enveloping Algebras and Geometric Representation Theory} (2012), Oberwolfach Reports, pp.~733--809, DOI \texttt{10.4171/owr/2012/13}. \url{https://doi.org/10.4171/owr/2012/13}. The record describes Oberwolfach Reports as documenting mathematical workshops and their open-problem discussions; no precise problem page or passage is supplied.
\item Supplied literature reference: \emph{Prime Representations from a Homological Perspective}. \url{https://arxiv.org/abs/1112.6376}.
\item Supplied literature reference: \emph{Prime Representations from a Homological Perspective}. \url{https://doi.org/10.1007/s00209-012-1088-7}.
\item Literature assessment dated exactly \texttt{2026-08-21}: the supplied curation labels the problem open and reports that the self-extension criterion is proved in type \(\mathfrak{sl}_2\) and for broad classes of prime representations, while a converse characterization for all simple quantum-affine representations remains conjectural. This is supplied dated curation derived from the OWR record, not an independently verified literature assessment. Its time-sensitive status and cited claims require independent verification.
\end{itemize}
\end{document}
